% ECONOMICS_PROBLEM: {"id": "C73-159", "title": "Strategic queue stability", "jel_code": "C73", "source_id": "OP-0286"}
% !TeX program = xelatex
\documentclass[11pt,a4paper]{article}
\usepackage[margin=23mm]{geometry}
\usepackage{fontspec}
\setmainfont{Latin Modern Roman}
\usepackage{amsmath,amssymb,mathtools}
\usepackage{xcolor,enumitem,booktabs,longtable,array,xurl,hyperref}
\definecolor{muted}{HTML}{53636C}
\hypersetup{colorlinks=true,urlcolor=blue,linkcolor=blue}
\setlength{\parindent}{0pt}
\setlength{\parskip}{5pt}
\newcommand{\R}{\mathbb R}
\newcommand{\E}{\mathbb E}
\newcommand{\N}{\mathbb N}
\newcommand{\ind}{\mathbf 1}
\newcommand{\Var}{\operatorname{Var}}
\newcommand{\Cov}{\operatorname{Cov}}
\newcommand{\supp}{\operatorname{supp}}
\DeclareMathOperator*{\argmin}{arg\,min}
\DeclareMathOperator*{\argmax}{arg\,max}
\newcommand{\entrymeta}[3]{{\sffamily\small\color{muted}\textbf{\texttt{#1}}\quad|\quad #2\par #3\par}}
\newcommand{\Problemref}[1]{\texttt{#1}}
\newcommand{\JELref}[2]{\texttt{#2} (JEL \texttt{#1})}
\begin{document}
\section*{Strategic queue stability}
\textbf{Problem ID:} \texttt{C73-159}\quad \textbf{JEL:} \texttt{C73}\par
% BEGIN ECONOMICS STATEMENT
\label{problem:OP-0286}
{\sffamily\small\color{muted}Original subject group: Game theory and strategic interaction\par}
\label{definitions:problem:30007568}
\textbf{Audit classification:} Published research agenda. Tardos Theorem 8 supplies a factor-three result for its own tiny-buffer model with additional \ensuremath{\lambda}\_i<1/2. The architecture-specific optimization below is editorial; neither factor is universal.
\subsection*{Definitions and precise research target}
\textbf{Model and research target.} For the fully specified strategic queue model and learner class, determine the least load-slack factor \(c^*\) defined below that guarantees a time-uniform bound on expected total backlog. Seek graph- and buffer-dependent constructive guarantees and matching destabilization examples. Source factors two and three are established benchmarks in their own models and do not certify a numerical gap for this extension.

There are \(n\) source queues and \(m\) servers on a fixed bipartite graph \(E\). Queue \(i\) receives independent Bernoulli arrivals of rate \(\lambda_i\in(0,1)\); server \(j\) processes at most one packet with independent Bernoulli success probability \(\mu_j\in(0,1]\) per round and has a finite packet buffer \(b_j\). Each nonempty source chooses a compatible server, and rejected packets remain at the source. For a finite initial state, strong stability means \(\sup_{t\ge0}\mathbb E[\sum_iQ_i(t)+\sum_jH_j(t)]<\infty\), where \(Q_i\) and \(H_j\) count source and server-buffer packets respectively. \emph{Definitions and maintained assumptions (editorial unless a source definition is named).} Specify the order of arrivals, routing, admission and service within a round; compatible edges \(E\); how an occupied buffer admits a new packet; collision priorities; and whether a packet can leave a source and be served in the same round. All packets must be counted exactly once in \(\sum_iQ_i+\sum_jH_j\). Central feasibility refers to strictly stabilizable arrival rates for these same rules, not only \(\sum_i\lambda_i<\sum_j\mu_j\). Let \(\Lambda^{\rm cent}\) denote that open stability region. To avoid multiplying Bernoulli success probabilities above one, define a load-slack factor
\[c^* = \inf\{c\ge1:\lambda^0/c\text{ is strongly stable for every }\lambda^0\in\Lambda^{\rm cent}\text{ and every }L\in\mathcal L\}.\]
Fix finite initial backlog and a high-probability window-regret bound \(R_i(s,w)\le r(w,\delta)\) with probability \(1-\delta\), together with its counterfactual routing payoff and filtration. Stability is the time-uniform first-moment bound in the statement; the constant may depend on the arrival vector’s slack. Global expected no-regret alone does not supply the required window event.
\subsection*{Variants and assumption changes}
\begin{enumerate}
\item \textbf{Tiny buffers (source baseline, editorial extension).} Begin with full compatibility, single-packet server buffers and the source’s exact learner/admission hypotheses. Retain its \(\lambda_i<1/2\) condition before seeking an improved slack factor.
\item \textbf{Restricted network (editorial).} Fix a noncomplete compatibility graph and capacities \(b_j\). Characterize stability through cut constraints and admissible learner feedback, and separate topological obstructions from failures caused by collision policies.
\end{enumerate}
\subsection*{References and source audit}
\begin{enumerate}
\item Packet routing as a dynamic game; Theorem8. Supports the stated research area; exact displayed specification is editorial. Theorem8 requires \(\lambda_i<1/2\) as well as factor-three total-load slack. \url{https://epubs.siam.org/doi/full/10.1137/25M1808064}
\end{enumerate}
\begin{enumerate}\raggedright
\item\hypertarget{definitions:source:30007568:1}{} Éva Tardos. 2026. \emph{Learning and the Price of Anarchy in Games}. Learning outcomes in dynamic games; Packet routing as a dynamic game; Theorem8.
\url{https://epubs.siam.org/doi/full/10.1137/25M1808064}.
DOI: \url{https://doi.org/10.1137/25M1808064}.
\end{enumerate}

\subsection*{Common conventions: Source mathematical conventions}

These conventions apply to each entry unless explicitly replaced.

\textbf{Models and identification.} A primitive model \(m\) specifies all probability laws, feasible choices, information, equilibrium and measurement rules used by its target. The admissible class \(\mathcal M\) is fixed independently of outcomes. If \(P_O(m)\) is its observable probability law and \(\tau(m)\) the target, its identified set at observable law \(P\) is
\[
 \mathcal I_\tau(P)=\{\tau(m):m\in\mathcal M,\ P_O(m)=P\}.
\]
Point identification means that this set is a singleton. An empty set signals incompatibility of data and assumptions. Coordinate bounds need not describe the joint set. Bounds are sharp only if every retained target is attainable or arbitrarily approachable by compatible models. Finite bounds require explicit restrictions on unobserved quantities; unrestricted confounding, hidden wealth, extrapolation or arbitrary equilibrium selection can yield uninformative sets.

\textbf{Causal interventions.} A potential outcome \(Y_i(p)\) is the outcome under a complete policy regime \(p\), including timing, financing, information and market spillovers. The target population and units are fixed before treatment. With interference, use the complete assignment vector or a justified exposure mapping; individual no-interference notation alone cannot identify an economy-wide intervention. A difference of mean potential outcomes does not require the cross-world joint distribution. Conditioning on post-treatment migration, survival, enrollment or employment changes the estimand and needs separate justification.

\textbf{Statistical statements.} An estimator, confidence set, forecast or test specifies a sampling law, model class and loss/coverage criterion. Uniform \((1-\alpha)\)-coverage means \(\inf_{m\in\mathcal M}\Pr_m\{\tau(m)\in C_n\}\geq1-\alpha\), or an explicitly asymptotic version. The whole parameter space trivially has coverage, so informativeness requires a stated width, risk or power objective. A forecasting improvement is relative to a named benchmark, information set and evaluation population.

\textbf{Equilibrium and optimization.} An equilibrium comprises feasible plans, optimality, beliefs where needed, budget constraints and all market/resource clearing. The primitives or a stated benchmark define each correspondence \(\mathcal E(p;m)\). Nonemptiness is assumed or proved, never inferred just from compact policy sets. A selected-equilibrium target uses a declared selection rule; a robust target quantifies over all permitted equilibria. Use suprema or infima unless attainment is established. An \(\varepsilon\)-optimal policy is feasible and within \(\varepsilon\) of the finite optimum value. Report an empty feasible set separately.

\textbf{Welfare and accounting.} Normative utility, interpersonal normalization, social weights, numeraire, discounting and the use of public receipts are primitives. Behavioral utility need not coincide with the welfare criterion. Household welfare includes ownership income and rents once; adding profits again to compensating variation generally double-counts them. A monetary equivalent is defined by an explicit transfer and price convention, with monotonicity and range conditions. Average effects, mechanism contributions, transfers and welfare losses are different targets. Infinite sums require convergence and dynamic feasible plans require terminal or no-Ponzi conditions.

\textbf{Source versions.} A working-paper date, revision date and journal publication year are distinguished. A DOI for a journal version is not automatically the DOI of an earlier manuscript. Locators refer to the stated version and printed pages where specified. A metadata-only check does not verify a claim about a full-text conclusion. Every entry's source subsection narrows the provenance claim accordingly.
% END ECONOMICS STATEMENT
\end{document}
